% Transcription — fonds Alexandre Grothendieck, shelfmark 73, batch 1 (pages 1-20).
% Established from the facsimile archives/batches/73/batch-01.pdf, prepared
% from a copy of 73.pdf fetched through the project's own relay
% (https://grothendieck-relay.onrender.com/source/73.pdf), not uploaded by the
% user; per the skill transcribe-grothendieck. The batch file carries no cover
% sheet: PDF page k is archivists' page k. The folder is twenty-eight pages;
% this is the first of two batches. Pages 21-28 are batch 2, transcribed in
% parallel by another pass; its notation (double-struck letters as \mathbf,
% products written \lambda_3\cdots\lambda_{n-1}, boxed formulas kept boxed)
% was read and followed.
%
% Pass: Opus 5.5 (claude-opus-5-5), 2026-09-26 — first pass, unchecked against the pages by a human.
%
% Numbering. The pencilled numbers bottom-left of these leaves are not the
% PDF count: page 1 is pencilled 3, and pages 2-20 are pencilled 9-27 (seven
% ahead; pencil 1-2 and 4-8 have no scan here). Batch 2 continues with 28-35.
% \page follows the PDF count, as for every other folder and as batch 2 does;
% the discrepancy is recorded here only.
%
% Dating and title. The dating is the folder's own, copied verbatim from the
% catalogue; the group « Géométrie et topologie combinatoire » (69-89,
% 1976-[vers 1986]) holds it, and since the folder has a date of its own the
% group's range is not added. Title from the catalogue, nothing bracketed.
%
% Pages skipped: 2, 7, 11 and 13, cream sheets carrying nothing but the
% show-through of the neighbouring leaf. Pages summarised: none. There is no
% letter in this batch either, by him or by anyone else — every written leaf
% is his working notes, one hand, black then blue-black ink; batch 2 reports
% none in its pages, so the « lettre (s.d.) » of the catalogue is not among
% the scanned leaves as far as the two passes can see. Page 8 carries, up the
% right margin, a short vertical list in his hand that is not mathematics (it
% reads as a list of first names); it is left out.
%
% His pagination: page 4 carries a circled « 21 » top right (isolated; it
% belongs to no run visible here). Pages 14-19 carry his circled 1 to 6, and
% page 20 an uncircled « 8 » (7 is not in the folder); batch 2 continues the
% same run with 9-16 on pages 21-28. Recorded once per page with a \note.
%
% What the batch is about. Plane polygons, in two registers, with one
% intruding piece.
%   1      The scheme of n-gons: for E an affine plane over S, Polg_n(E)(S')
%          = families of vertices s_i in E(S') and lines d_i in Dr(E_S')
%          with s_i, s_{i+1} on d_i; an open U; sketches of degenerations
%          (consecutive sides on one line, a vertex doubled, a triangle as a
%          degenerate hexagon).
%   3-6    Polygons given by side vectors u_i with sum 0 and a relation
%          u = lambda_i u_{i-1} - lambda_{i+1} u_{i+1}; M_i = L_{i-1} - L_i;
%          p. 4 solves the recurrence u_1 -> u_n and counts dimensions
%          (n-2 in P^{n-1}, codim 3); p. 5 the quadrilaterals: M_4 -> P^2 via
%          mu_i = u_{i-1} ^ u_i, mu_0 + mu_2 = mu_1 + mu_3, mono, M_4 ≃ P^2
%          (over Z), explicit u_2, u_3 in the basis u_0, u_1; p. 6 the three
%          degenerate quadrilaterals (0,0,1,1), (1,0,0,1), (1,1,0,0) form a
%          projective frame, so M(Pi) is identified with the projective
%          envelope of A, and sketches of quadrilaterals.
%   8-10   Not polygons: « configurations » C, « opérations » Omega with
%          results f_omega : C -> V_omega, the « opérations effectuées »
%          Omega~ = coprod V_omega -> Omega, residual sets C(omega~), tests
%          T_nu (image of C x Omega^nu in Omega~^nu), and successive choices
%          f_i : F~_{i-1} -> P(Omega) building F~_i inside T_i — a sequential
%          search / test strategy « de longueur nu ». Written in reverse
%          order: p. 10 sets up the definitions, pp. 9 and 8 build on them.
%   12     Hypotheses for the general n-gon (u = lambda_i u_{i-1} -
%          lambda_{i+1} u_{i+1}, mu_i = u_{i-1} ^ u_i, alpha_i = u ^ u_i),
%          the relations among lambda, mu, alpha, the cyclic n-gons CI_n^*
%          mod Aff mapping to P^{n-1} x P^{n-1}.
%   14-19  His 1-6: the differentials df_i of the areas f_i = x_{i-1} y_i -
%          y_{i-1} x_i, the closed formula (y_{i-1} y_i) lambda_i =
%          (-lambda_0 y_{-1})(y_1+...+y_{i-1}) + (lambda_1 y_1)(y_0+...+y_{i-1})
%          and its vector form lambda_i u_{i-1} u_i = -lambda_0 u_{-1} u_{1i}
%          + lambda_1 u_1 u_{0i} with u_{ab} = s_b - s_a; the quadrilateral
%          (lambda_0 + lambda_1 = 0 at a non-flat vertex), an explicit
%          quadrilateral with vertices on two crossing lines, the 9-gon with
%          lambda_0 = 0, and the pentagon (coordinates in terms of the
%          f_i, a quadratic equation for x and its discriminant).
%   20     His 8: the four-term relation lambda_{i-1} u_{i-2} - lambda_i
%          (u_{i-1} + u_i) + lambda_{i+1} u_{i+1} = 0 (boxed), u_i in the
%          basis u_0, u_1, u_2 — the starting point of batch 2.
%
% Notation fixed once for the batch: double-struck letters as \mathbf; his
% script P for the set of subsets as \mathfrak{P}; his « ∐ » for the sum
% kept; the wedge u ^ v kept as \wedge where he writes it, and the plain
% juxtaposition « u_{i-1} u_i » (pp. 4, 15, 18) kept as written, since it
% denotes the same determinant; boxed formulas kept boxed. Struck pieces
% inside formulas are not set with \struck (which stays out of math) but
% omitted and recorded in a \note.
%
% Hardest pages: 4 (the prose in the middle, almost entirely lost), 5 (the
% closing paragraph), 8-10 (the vocabulary of the « opérations » and the
% interlinear additions), 12 (the right-hand hypotheses block, much
% overwritten) and 19 (the boxed cross terms of the discriminant, whose signs
% are overwritten). Readings to check first: « Polg » on p. 1; « Calcul
% intrinsèque » on p. 4; « enveloppe » on p. 6; « crible » on p. 9;
% « opérations » and « résiduel » on p. 10; « aplati » on p. 16; the value
% v written over a 4 on p. 17.

\documentclass[11pt,a4paper]{article}
% Le préambule commun à toutes les transcriptions du fonds Grothendieck.
%
% Il tient en une idée : **l'incertitude se compose, elle ne se résout pas.**
% Une transcription qui lisse un mot douteux a détruit la seule chose pour
% laquelle on la faisait — un lecteur doit pouvoir savoir, à chaque endroit,
% ce qui était sur le feuillet et ce qui a été ajouté. Les six macros qui
% suivent sont donc l'appareil critique, et non de la décoration.
%
% Les mêmes commandes sont reconnues par `scripts/render.mjs`, qui produit la
% vue de lecture du site. Une macro ajoutée ici doit l'être aussi là-bas,
% faute de quoi le rendu échouera bruyamment — ce qui est voulu : un rendu
% silencieusement faux est pire qu'un rendu absent.

\ProvidesPackage{grothendieck}[2026/08/08 Transcriptions du fonds Grothendieck]

\usepackage{amsmath,amssymb,amsthm}
\usepackage{mathrsfs}
\usepackage{stmaryrd}
% Les diagrammes commutatifs sont partout dans ce fonds : c'est la langue
% même de la géométrie algébrique de Grothendieck, et une transcription qui
% les rendrait en prose ne transcrirait rien.
\usepackage{tikz-cd}
% Français par défaut — c'est la langue des feuillets — mais l'anglais est
% chargé aussi : la traduction et le résumé sont des documents anglais, et la
% typographie française (espaces avant les deux-points) y serait une faute.
% Ces documents basculent par \selectlanguage{english} en tête de corps.
\usepackage[english,french]{babel}
\usepackage{fontspec}
\usepackage{geometry}
\geometry{a4paper,margin=25mm,marginparwidth=28mm}
\usepackage{marginnote}
\usepackage{xcolor}
% ulem, not soul. `soul`'s \st and \ul are built on a character-by-character
% scan that XeLaTeX's font machinery breaks: the document fails with an
% undefined control sequence rather than degrading. ulem does the same two jobs
% and is Unicode-engine safe. `normalem` stops it hijacking \emph, which the
% transcriptions use for Grothendieck's own emphasis.
\usepackage[normalem]{ulem}
\usepackage{hyperref}
\hypersetup{colorlinks=true,linkcolor=black,urlcolor=black,pdfborder={0 0 0}}

% --- Métadonnées du lot -----------------------------------------------------
% Reprises telles quelles par le script de rendu, qui les affiche en tête de
% la vue de lecture. Elles fixent ce que le fichier prétend couvrir : sans
% elles, une transcription flotte sans point d'attache dans un fonds de seize
% mille feuillets.

\newcommand{\folder}[1]{\gdef\grothFolder{#1}}
\newcommand{\batch}[1]{\gdef\grothBatch{#1}}
\newcommand{\leaves}[2]{\gdef\grothFirst{#1}\gdef\grothLast{#2}}
\newcommand{\foldertitle}[1]{\gdef\grothFoldertitle{#1}}
\gdef\grothFolder{?}\gdef\grothBatch{?}\gdef\grothFirst{?}\gdef\grothLast{?}\gdef\grothFoldertitle{}

% --- Le résumé --------------------------------------------------------------
%
% Il ouvre la lecture modernisée, sous le titre « Résumé », et il est composé
% en petit. Ce n'est pas un ornement : c'est le seul passage du document qui
% ne soit pas des mathématiques, et un lecteur doit pouvoir le reconnaître
% avant de l'avoir lu — puis le sauter s'il connaît déjà le sujet, ou s'y
% arrêter s'il ne le connaît pas. Le filet qui le suit marque la reprise du
% corps.

\newenvironment{resume}
  {\small\setlength{\parskip}{0.45em}}
  {\par\medskip\hrule\medskip}

% --- Mots-clés ---------------------------------------------------------------
%
% En anglais, en clôture du résumé de la lecture modernisée : le vocabulaire
% moderne sous lequel chercher ce que les feuillets construisent. Le manifeste
% du site (`npm run manifest`) les extrait pour étiqueter la cote entière —
% cette ligne est la source unique des tags, il n'y a pas de fichier à côté.
\newcommand{\keywords}[1]{\par\smallskip\noindent
  {\footnotesize\textsc{Keywords} --- \textit{#1}}\par}

% --- L'appareil critique ----------------------------------------------------

% Le numéro de feuillet, en marge. C'est l'ancre qui permet de régler un
% désaccord sur une formule en regardant *un* feuillet plutôt qu'une chemise
% de six cents. Le site s'en sert aussi pour tourner les pages du fac-similé
% quand on fait défiler la transcription.
\newcommand{\leaf}[1]{%
  \marginnote{\footnotesize\color{gray!70}\textsf{#1}}%
  \label{leaf:#1}\ignorespaces}

% Illisible : on ne devine pas. Un mot inventé qui se lit comme les autres est
% le pire résultat possible.
\newcommand{\ill}{\textcolor{red!60!black}{[\,\ldots\,]}}

% Lecture douteuse : le mot est donné, et signalé comme incertain.
\newcommand{\uncertain}[1]{\uline{#1}}

% Ajout de l'éditeur — une lettre manquante, un mot sous-entendu.
\newcommand{\add}[1]{\textcolor{blue!50!black}{[#1]}}

% Biffé par Grothendieck lui-même. Ce qu'il a rayé fait partie du document :
% une rature dit souvent où la pensée a changé de direction.
\newcommand{\struck}[1]{\sout{#1}}

% Note du transcripteur, sur l'état matériel du feuillet.
\newcommand{\note}[1]{\par\smallskip{\small\color{gray!60!black}\textit{[#1]}}\par\smallskip}

% Note marginale de Grothendieck, distincte du corps du texte.
\newcommand{\marginal}[1]{\par\smallskip\noindent\hspace{1em}%
  {\small\color{gray!60!black}#1}\par\smallskip}

% --- Titre ------------------------------------------------------------------

\AtBeginDocument{%
  \begin{center}
    {\large\bfseries Fonds Alexandre Grothendieck --- cote n\textsuperscript{o}~\grothFolder}\\[2pt]
    {\itshape\grothFoldertitle}\\[4pt]
    {\small feuillets \grothFirst--\grothLast\ (lot \grothBatch)}\\[2pt]
    {\footnotesize\color{gray!60!black}Transcription établie d'après le fac-similé numérisé
     par l'Université de Montpellier.\\
     Les lectures incertaines sont soulignées, les illisibles notées [\,\ldots\,],
     les ajouts entre crochets.}
  \end{center}
  \vspace{1em}}

\endinput


\folder{73}
\batch{1}
\pages{1}{20}
\dating{[à partir de 1976]}
\watermark{Édition de démonstration}
\foldertitle{Polygones : notes manuscrites (s.d.), lettre (s.d.).}

\begin{document}

\page{1}
$E$ plan affine$/S$ \qquad $n$ entier $\geqslant 3$

\[
\mathrm{Polg}_n(E)(S') = \Bigl\lbrace s_0, s_1, \ldots, s_{n-1}, d_0, d_1,
\ldots, d_{n-1} \Bigm|
\begin{array}{l}
s_i \in E(S') \simeq \Gamma(E_{S'}/S') \\
d_i \in \mathrm{Dr}(E_{S'}) \quad (0 \leqslant i \leqslant n-1)
\end{array} \Bigr.
\]
\[
\Bigl. s_0, s_1 \prec d_0 \,;\ s_1, s_2 \prec d_1 \,;\ \ldots \,;\
s_{n-1}, s_0 \in d_{n-1} \Bigr\rbrace
\]
\note{le « P » initial est doublé, comme une lettre ajourée ; le nom
« Polg » est \uncertain{lu}. « $\mathrm{Dr}$ » : les droites. Le signe
$\prec$ (incidence) est employé pour les premiers couples, $\in$ pour le
dernier}

\[
\mathcal{U}_{\underline{S}_0} = \Bigl\lbrace \bigl( (P_{\underline{s}})_{\underline{s}
\in \underline{S}}, (d_{\underline{a}})_{\underline{a} \in \underline{A}} \bigr)
\in \mathbf{P}\mathrm{ol}_n \Bigm| \ldots \ \forall\, s \in \underline{S}
\Bigr\rbrace
\]
\note{la condition, à droite du trait, est biffée (on lit encore
« $P_s \notin \ldots d_{s-2}$ ») ; au-dessus, « $d_{\underline{a}}$ »
seul, non biffé ; « $\forall\, s \in \underline{S}$ » est lui aussi
barré d'un trait. L'indice de $\mathcal{U}$ est \uncertain{lu}}

$s_2 \in d_0$ \qquad $s_0, d_0, s_1, d_1, \ldots, s_{n-3}, d_{n-3}, s_{n-2}$

\note{en marge gauche, une figure : les droites $d_0$, $d_1$, $d_{n-2}$,
$d_{n-1}$ et les sommets $s_0$, $s_1$, $s_2$, $s_{n-2}$, $s_{n-1}$ du
polygone. Plus bas, trois croquis de dégénérescences : $s_0$, $s_1$,
$s_2 = s_3$ alignés sur une même droite ($d_0$, $d_1$, avec en dessous
« \uncertain{$u\,d_{12}$} ») ; deux droites sécantes portant deux points
cerclés ; un triangle vu comme hexagone dégénéré, de sommets $s_0 = s_1$,
$s_2 = s_3$, $s_4 = s_5$ et de côtés $d_0 = d_1$, $d_2 = d_3$,
$d_4 = d_5$}

\[
\underbrace{d_0 = d_1 = \cdots = d_i}_{s_0, s_1, \ldots, s_{i+1}} \neq d_{i+1}
\ni s_{i+2}
\]
\note{les appartenances sont écrites verticalement ($\Downarrow$ sous
$d_i$ et $d_{i+1}$) ; figure : $s_0$, $s_i$, $s_{i+1} = s_{i+2}$ sur
$d_0 = \cdots = d_i$, puis $d_{i+1}$ qui repart du point double}

\page{3}
\[
L_0(\lambda)\,\sigma = 0 \qquad L_1(\lambda)\,\sigma = 0
\qquad\qquad L_0(\lambda) = L_1(\lambda) \qquad L(\lambda)\,\sigma =
\]
\note{une troisième ligne de la colonne de gauche est biffée}

\[
u\,\sigma = 0 \qquad\qquad \Sigma \qquad\qquad
\underbrace{L_0 = L_1 = \cdots = L_{n-2}}_{u} \quad L_{n-1}
\]

\[
\underbrace{(L_{n-2} - L_{n-1})}_{M_{n-1}}\, u_{n-1} = u\,.\,\sigma
\]

\[
\begin{aligned}
M_1 u_1 &= (L_0 - L_1)\, u_1 = \lambda_0 u_{-1} u_{12} - \lambda_1 u_1 u_{02}
 + \lambda_2 u_1 u_2 = \\
M_2 u_2 &= (L_1 - L_2)\, u_2 = (\lambda_0 u_{-1} u_{13} - \lambda_1 u_1 u_{03})
 + \lambda_3 u_2 u_3 \quad \mathrm{mod}
\end{aligned}
\]
\note{le signe entre $\lambda_0 u_{-1} u_{12}$ et $\lambda_1 u_1 u_{02}$
est surchargé ; « mod » reste en suspens, et le reste de la page est vide}

\page{4}
\note{numéro « 21 » cerclé en haut à droite, isolé}
\[
\lambda_2 u_1 u_2 = -\lambda_0 u_{-1} u_1 + \lambda_1 u_1 (u_0 + u_1)
\]
\note{les deux $u_1$ de droite sont barrés d'un trait oblique}
\[
\underbrace{\lambda_n}_{\lambda_0} u_{n-1} \underbrace{u_n}_{u_0}
= -\lambda_0 \underbrace{u_{-1}}_{u_{n-1}} \underbrace{u_{1n}}_{u_{0n} - u_0}
+ \lambda_1 u_1 u_{0n}
\]
\ill{} \ill{} $\lambda_*$

la dim.\ de la variété des $\mu_*$ tels que --- \ill{} $n-2$
($\subset \mathbf{P}^{n-1}$)

\emph{dans} le cas \ill{} $\sum u_i = 0$

\ill{} $n-4$ \ill{} la condition $\sum u_i = 0$ \ill{} $\mathbf{P}^{n-1}$
\ill{} \ill{} codim.\ 3 dans \ill{}

\marginal{$(c_1, \ldots, c_{n-2})$ ; $c_1 + c_3 + \cdots + c_{n-3} = 0$ ;
$c_2 + c_4 + \cdots + c_{n-2} = 0$}
\note{ces lignes de prose sont écrites très vite ; seuls les symboles se
lisent sûrement. Un trait sépare ensuite la page en deux}

(\uncertain{Calcul intrinsèque})

\note{en tête de la seconde moitié, quelques essais biffés
(« $\lambda_0$ », « $u$ -- $u_1, u_1$ », « $-\lambda_2 u_2$ ») ;
« $u_0, u_1$ » souligné d'une accolade, et encadré : « $u_0, u_1 ; u$ »}

\[
\begin{array}{lll}
u_1 & \lambda_1 u_0 - \lambda_2 u_2 = u & \Rightarrow u_2 \\
u_2 & \lambda_2 u_1 - \lambda_3 u_3 = u & \Rightarrow u_3 \\
u_{n-2} & \lambda_{n-2} u_{n-3} - \lambda_{n-1} u_{n-1} = u & \Rightarrow u_{n-1} \\
u_{n-1} & \lambda_{n-1} u_{n-2} - \lambda_n u_n = u & \Rightarrow u_n = u_0 \\
u_n & \lambda_n u_{n-1} - \lambda_{n+1} u_{n+1} = u & \Rightarrow u_{n+1} = u_1
\end{array}
\]
En fonction de $u_0, u_1, u$ ; les deux dernières lignes donnent chacune une
relation entre $u_0, u_1, u, \lambda_*$.
\note{« en fonction de $u_0, u_1, u$ » est écrit à droite d'un trait
vertical qui embrasse les trois premières lignes ; « relation entre
$u_0, u_1, u, \lambda_*$ » est écrit deux fois, en face des deux dernières}

\page{5}
\note{figure : un quadrilatère $s_0 s_1 s_2 s_3$, de côtés orientés
$u_0, u_1, u_2, u_3$}
\[
\mu_0 + \mu_2 = \mu_1 + \mu_3 \qquad\qquad \mu_3 = \mu_0 - \mu_1 + \mu_2
\]
\note{une ligne biffée sous la première formule}

\[
\mathbf{M}_4 \longrightarrow \mathbf{P}^2 \qquad
(u_0, u_1, u_2, u_3) \longmapsto (\underbrace{u_3 \wedge u_0}_{\mu_0},
\underbrace{u_0 \wedge u_1}_{\mu_1}, \underbrace{u_1 \wedge u_2}_{\mu_2})
\]
\note{le $\mathbf{M}$ est une capitale ajourée, qu'on rend par un gras}

\[
\begin{aligned}
u_0 \wedge u_2 &= u_0 \wedge (u_0 + u_1) \\
 &= u_0 \wedge (u_1 + u_3) \\
 &= -\underbrace{u_0 \wedge u_1}_{\mu_1} + \underbrace{u_3 \wedge u_0}_{\mu_0}
\end{aligned}
\]
\note{il faut lire $-u_0 \wedge (u_1 + u_3)$ à la deuxième ligne ; le
signe est surchargé et la première ligne reste telle qu'elle est écrite}

\[
\begin{aligned}
u_0 \wedge u_2 &= \mu_0 - \mu_1 \\
u_1 \wedge u_3 &= \mu_0 - \mu_2 \\
u_2 \wedge u_0 &= \mu_1 - \mu_3 = \mu_1 - \mu_0 = -u_0 \wedge u_2 \\
u_3 \wedge u_1 &= \mu_3 - \mu_0 = \mu_2 - \mu_1 = -u_1 \wedge u_3
\end{aligned}
\]
\note{au-dessus, « $u_0 \wedge u_2$ » biffé}

$\mathbf{M}_4 \to \mathbf{P}^2_{\mathbf{Z}}$ \struck{injectif} mono (car la
\uncertain{connaissance} des $\mu_0, \mu_1, \mu_2$ \uncertain{indique} celle
de $\mu_3$, \uncertain{donc} les $u_i \wedge u_j$ \ldots)

épi ?

$\mu_0, \mu_1, \mu_2$, \struck{$\mu_3$} \ill{} donnés,
$\mu_0 + \mu_2 = \mu_1 + \mu_3$, $\mu_3$ \ill{} \ill{}

\struck{\ill{}} $\mu_1 \neq 0$. On prend $u_0, u_1$ comme base (de $V$),
$u_0 \wedge u_1$ base de $\bigwedge^2 V$,
\[
u_2 = \alpha u_0 + \beta u_1, \qquad u_3 = -(u_0 + u_1 + u_2)
\]

\[
\begin{array}{ll}
u_0 \wedge u_1 = 1\,.\,e & \mu'_1 = 1 \\
u_1 \wedge u_2 = -\alpha\, e & \mu'_2 = -\alpha \\
u_2 \wedge u_3 = (\beta - \alpha)\, e & \mu'_3 = \beta - \alpha \\
u_3 \wedge u_0 = (1 + \beta)\, e & \mu'_0 = 1 + \beta
\end{array}
\qquad
\begin{aligned}
u_2 \wedge u_3 &= -u_2 \wedge u_0 - u_2 \wedge u_1 \\
 &= \beta \underbrace{u_0 \wedge u_1}_{e} - \alpha\, e \\
u_3 \wedge u_0 &= -u_1 \wedge u_0 - u_2 \wedge u_0 \\
 &= e + \beta\, e
\end{aligned}
\]
\note{la ligne $u_2 \wedge u_3$ et la valeur $\mu'_3$ sont surchargées
(on lit $\beta - \alpha$ sous la rature) ; un trait sépare
$\mu'_3$ de $\mu'_0$}

\[
\alpha = -\frac{\mu_2}{\mu_1} \qquad \beta - \alpha = \frac{\mu_3}{\mu_1}
\qquad \text{i.e.} \qquad \beta = \frac{\mu_3}{\mu_1} + \alpha
= \frac{\mu_3 - \mu_2}{\mu_1}
\]
\[
u_2 = -\frac{\mu_2}{\mu_1} u_0 + \frac{\mu_3 - \mu_2}{\mu_1} u_1
\qquad\qquad \mathbf{M}_4 \simeq \mathbf{P}^2_{\mathbf{Z}}
\]
\[
u_3 = \Bigl( \frac{\mu_2}{\mu_1} - 1 \Bigr) u_0
+ \underbrace{\frac{\mu_2 - \mu_3 - \mu_1}{\mu_1}}_{-\mu_0} u_1
\qquad\qquad
u_2 = 0 \iff \boxed{\mu_2 = 0,\ \mu_3 = 0}\,,\quad \mu_0 = \mu_1
\]
\note{après $\iff$, « $\mu_1 \neq 0$ » biffé au-dessus de la condition
encadrée}

Les \uncertain{quadrilatères} \uncertain{dégénérés}, \ill{} un côté nul
\add{(}$a_i$ ($i = 0, 1, 2, 3$)\add{)}, correspondent aux points \ill{}
\ill{} de $\mathbf{P}^2$ ($\mu_0 + \mu_2 = \mu_1 + \mu_3$) $\subset
\mathbf{P}^3$ $(\mu_0, \mu_1, \mu_2, \mu_3)$ correspondant à deux
coordonnées \ill{} \ill{} (\ill{} \ill{} \ill{}) \ill{} \ill{} \ill{}
\ill{} \ill{}
\note{« $a_i$ ($i = 0, 1, 2, 3$) » est ajouté dans l'interligne ; un mot
biffé avant « $\mu_0 + \mu_2$ ». La phrase continue en tête de la page 6}

\page{6}
\ill{} alignés (grâce \uncertain{aux} vecteurs).
$(0, 0, 1, 1)$ $(1, 0, 0, 1)$ $(1, 1, 0, 0)$ ne \uncertain{sont} \ill{}
\uncertain{coplanaires}, i.e.\ la matrice
\[
\begin{pmatrix} 0 & 0 & 1 & 1 \\ 1 & 0 & 0 & 1 \\ 1 & 1 & 0 & 0 \end{pmatrix}
\]
de \uncertain{rang} 3 \ill{})

Donc ils forment une base projective de $\mathbf{P}^2$, donc,
\struck{\ill{}} si $\Pi = (S, A)$ est un quadrilatère combinatoire,
$\mathbf{M}(\Pi)$ s'identifie à l'\uncertain{enveloppe} projective de $A$.

\note{sous le texte, cinq croquis de quadrilatères : un quadrilatère
ouvert, un quadrilatère croisé en « Z », deux droites sécantes portant
quatre points, un trapèze avec ses diagonales, un quadrilatère aplati avec
ses diagonales}

\page{8}
\[
\mathcal{F}_\nu \subset T_\nu \subset \widetilde{\Omega}^\nu \qquad
T_\nu = \mathrm{Im}(C \times \Omega^\nu \to \widetilde{\Omega}^\nu)
\]
\note{au-dessus, « $\mathcal{F}_\nu \subset \widetilde{\Omega}^\nu$ »
biffé. L'égalité $T_\nu = \mathrm{Im}(\ldots)$ est écrite verticalement
(« $\shortparallel$ ») sous $T_\nu$. Les notations sont celles de la
page 10 ; les pages 8--10 sont écrites dans l'ordre inverse de la
lecture}

\[
\mathcal{F}_i \subset T_i \subset \widetilde{\Omega}^i \quad (1 \leqslant i
\leqslant \nu) \qquad\qquad
\mathrm{pr}_{i-1,i}(\mathcal{F}_i) \subset \mathcal{F}_{i-1} \quad
(2 \leqslant i \leqslant \nu)
\]

\[
\mathcal{F}_i = \bigl( (\widetilde{\omega}_1, \ldots, \widetilde{\omega}_i)
\bigm| (\widetilde{\omega}_1, \ldots, \widetilde{\omega}_{i-1}) \in
\mathcal{F}_{i-1},\ \omega_i \ldots
\qquad \widetilde{\Omega}^{i-1} \times \Omega
\]
\[
\mathcal{F}_i = T_i \cap (\widetilde{\Omega}^{i-1} \times \widetilde{\Omega}
\,|\, \mathcal{F}_{i-1} \qquad T_{i-1} \times \Omega
\]

\[
\begin{array}{l}
f_1 \subset \Omega \\
f_2 \subset \Omega \times \Omega \\
\quad \cdots \\
f_i \subset \widetilde{\Omega}^{i-1} \times \Omega, \quad
 f_i \subset T_{i-1} \times \Omega \\
\quad \cdots \\
f_\nu \subset \widetilde{\Omega}^{\nu-1} \times \Omega
\end{array}
\]

\begin{tikzcd}[column sep=small, row sep=small]
\widetilde{\Omega}^{i-1} \times \Omega \supset f_i \\
\widetilde{\Omega}^i \quad \mathcal{F}_i \arrow[u] \\
\widetilde{\Omega}^i \times \Omega \supset f_{i+1} \arrow[u, "\mathrm{pr}_1"]
\end{tikzcd}

\note{la flèche du haut porte une marque raturée ; celle du bas
« $\mathrm{pr}_1$ » (\uncertain{lu})}

\[
\mathcal{F}_i \colon \quad
\mathcal{F}_i = \Bigl\lbrace \bigl( \widetilde{\omega}_1(x), \ldots,
\widetilde{\omega}_{i-1}(x), \omega_i(x) \bigr) \Bigm| x \in C,\
\bigl( \omega_1(x), \ldots, \omega_{i-1}(x), \omega_i \bigr) \in f_i
\Bigr\rbrace
\subset \widetilde{\Omega}^i
\]
\note{au-dessus, « $f_i \subset T_{i-1} \times \Omega$ » biffé ; sous
la condition, en indice, « $\omega_1, \ldots, \omega_i \in \Omega^i$ »}

\[
\begin{aligned}
f_i &\subset T_{i-1} \times \Omega \subset \widetilde{\Omega}^{i-1}
 \times \Omega \\
\mathcal{F}_i = \widetilde{f}_i &\subset T_i \subset
 \widetilde{\Omega}^{i-1} \times \widetilde{\Omega} = \widetilde{\Omega}^i \\
 &= T_i \cap (\widetilde{\Omega}^i \,|\, f_i)
\end{aligned}
\]
\note{une flèche relie les deux premières lignes, avec un mot \ill{} à
gauche}

\page{9}
\[
\boxed{\mathcal{F}_1 \subset \Omega} \quad \text{i.e.} \quad
\mathcal{F}_1 \in \mathfrak{P}^*(\Omega)
\]
\[
\widetilde{\mathcal{F}}_1 = (\widetilde{\Omega} \,|\, \mathcal{F}_1) \cap T_1
\subset \widetilde{\Omega} = \widetilde{\Omega}_1
\]
\[
\boxed{f_2 \colon \widetilde{\mathcal{F}}_1 \longrightarrow
\mathfrak{P}^*(\Omega)} \qquad\qquad \widetilde{\Omega} \times \Omega
\]
\note{il écrit l'ensemble des parties par un « P » gothique, ici avec un
astérisque : on le rend par $\mathfrak{P}$}

$(\widetilde{\omega}_1, \omega_2) \in \widetilde{\Omega} \times \Omega$
\uncertain{compatible} avec $(\mathcal{F}_1, f_2)$
\begin{itemize}
  \item[] si a) $\widetilde{\omega}_1 \in \widetilde{\mathcal{F}}_1$, i.e.\
  $\omega_1 \in \mathcal{F}_1$ et $\widetilde{\omega}_1 \in T_1$ ;
  \item[] b) $\omega_2 \in f_2(\widetilde{\omega}_1)$.
\end{itemize}
\note{un symbole biffé devant la première ligne, et une ligne biffée
entre a) et b)}

\[
\underbrace{\widetilde{\mathcal{F}}_2 \subset \widetilde{\Omega}_2 =
\widetilde{\Omega} \times \widetilde{\Omega}} = \bigl\lbrace
(\widetilde{\omega}_1, \widetilde{\omega}_2) \bigm| \omega_1 \in
\mathcal{F}_1,\ (\widetilde{\omega}_1 \times \widetilde{\omega}_2) \in T_2,\
\widetilde{\omega}_2 \in f_2(\widetilde{\omega}_1) \bigr\rbrace
\]
\[
f_3 \colon \widetilde{\mathcal{F}}_2 \longrightarrow \mathfrak{P}(\Omega)
\]
\[
f_\nu \colon \widetilde{\mathcal{F}}_{\nu-1} \longrightarrow
\mathfrak{P}(\Omega) \quad \text{donc} \qquad
\widetilde{\mathcal{F}}_\nu \subset \widetilde{\Omega}^\nu \quad
(\text{en fait, } \widetilde{\mathcal{F}}_\nu \subset T_\nu)
\]
\note{sous $\widetilde{\mathcal{F}}_{\nu-1}$, « $\cap\,
\widetilde{\Omega}^{\nu-1}$ » écrit verticalement}

$(\mathcal{F}_1, \ldots, \mathcal{F}_\nu)$ \uncertain{s'appelle} une
\struck{\ill{} \ill{}} \emph{\uncertain{crible} de longueur $\nu$}

\textbf{NB} certains $\widetilde{\mathcal{F}}_i$ peuvent être $\emptyset$.
\note{« crible de longueur $\nu$ » est ajouté au-dessus du mot biffé et
souligné}

\[
\begin{array}{ll}
\mathcal{F}_1 = f_1 \subset \Omega & \widetilde{\Omega} \times
\widetilde{\Omega} \cdots \times \widetilde{\Omega} \\
f_2 \subset \widetilde{\Omega} \times \Omega & \mathcal{F} \\
f_3 \subset \widetilde{\Omega} \times \widetilde{\Omega} \times \Omega & \\
f_\nu \subset \underbrace{\widetilde{\Omega} \times \widetilde{\Omega}
\cdots \widetilde{\Omega}}_{\nu-1} \times \Omega &
\end{array}
\]

\page{10}
$\mathcal{C}$ « configurations »

$\Omega$ « \uncertain{opérations} » (virtuelles)

$\omega \in \Omega$ \quad $\omega \longmapsto (f_\omega \colon \mathcal{C}
\to V_\omega)$ \qquad on pose $f_\omega(x) = \omega(x)$

$V_\omega$ = « ens.\ des résultats (virtuels) de l'opération $\omega$ »
\note{« virtuels » est ajouté dans l'interligne ; l'égalité
$V_\omega = \ldots$ est écrite verticalement sous $V_\omega$}

$\widetilde{\Omega}$ « opérations effectuées » :
\[
\widetilde{\Omega} = \coprod_{\omega \in \Omega} V_\omega
\xrightarrow{\ \pi\ } \Omega
\]

\[
C(\widetilde{\omega}, x) = f_\omega^{-1}\bigl(\lbrace f_\omega(x)
\rbrace\bigr) = C(\omega(x))
\]
ensemble des \uncertain{la conf.} $x \in C$ \ill{} par l'opér.\ $\omega$.

De \uncertain{même} on définit
\[
C(\omega_1, \ldots, \omega_\nu; x) = \bigcap_{1 \leqslant i \leqslant \nu}
C(\omega_i; x)
\]
ensemble des $x \in C$ \ill{} par la suite d'op.\ $\omega_1, \ldots,
\omega_\nu$.
\note{ces deux définitions sont embrassées à gauche par une accolade,
qu'une longue flèche relie au bas de la page}

\[
\widetilde{\Omega} \longrightarrow \mathfrak{P}(\mathcal{C}), \quad
\widetilde{\omega} \mapsto C(\widetilde{\omega}) \qquad\qquad
\widetilde{\omega} = (\omega, v) \longmapsto f_\omega^{-1}(\lbrace v \rbrace)
\]
\textbf{NB} \ill{} \ill{} \ill{} pas les \ill{} \ill{}
\note{« NB » souligné, suivi d'un mot souligné, \ill{}}

$\widetilde{\omega} \longmapsto C(\widetilde{\omega})$ \quad ensemble
\uncertain{résiduel} pour l'op.\ eff.\ $\widetilde{\omega}$

\[
\widetilde{\Omega}^\nu \longrightarrow \mathfrak{P}(C), \qquad
(\widetilde{\omega}_1, \ldots, \widetilde{\omega}_\nu) \longmapsto
C(\widetilde{\omega}_1, \ldots, \widetilde{\omega}_\nu) =
\bigcap_{1 \leqslant i \leqslant \nu} C(\widetilde{\omega}_i)
\]
ensemble \uncertain{résiduel} pour la suite d'opérations effectuées
$\widetilde{\omega}_1, \ldots, \widetilde{\omega}_\nu$.

$T_\nu \subset \widetilde{\Omega}^\nu$ \quad \emph{tests} \uncertain{virtuels}
de longueur $\nu$ / effectués de longueur $\nu$
\note{« $\Omega^\nu$ » écrit au-dessus de $\widetilde{\Omega}^\nu$ ;
« effectués de longueur $\nu$ » est une seconde ligne, ajoutée sous la
première}

\[
C \times \Omega \longrightarrow \widetilde{\Omega}, \quad (x, \omega)
\longmapsto \omega(x) \qquad \Longrightarrow \qquad
C \longrightarrow \Gamma(\widetilde{\Omega}/\Omega)
\]
\[
C \times \Omega^\nu \longrightarrow \widetilde{\Omega}^\nu
\]
\[
C \longrightarrow \Gamma(\widetilde{\Omega}^\nu/\Omega^\nu), \qquad
x \longmapsto \bigl( (\omega_1, \ldots, \omega_\nu) \longmapsto
(\omega_1(x), \ldots, \omega_\nu(x)) \bigr)
\]
\note{la flèche $\Longrightarrow$ remplace une flèche biffée}

L'image de $C \times \Omega^\nu$ \uncertain{s'appelle} ens.\ des tests
\uncertain{effectués} de long.\ $\nu$ ; \uncertain{soit} $T_\nu$ l'ens.\
de ceux-ci.

\page{12}
\note{figure : quatre côtés consécutifs $u_{i-1}, u_i, u_{i+1}, u_{i+2}$
d'un polygone, sommets $s_{i-1}, \ldots, s_{i+3}$}

hyp.\ \uncertain{données} :
\[
\left\lbrace
\begin{aligned}
& u = \lambda_i u_{i-1} - \lambda_{i+1} u_{i+1} \qquad i = 0, 1, \ldots, n-1 \\
& \textstyle\sum u_i = 0
\end{aligned}
\right.
\]
(déf.) :
\[
\left\lbrace
\begin{aligned}
& u_{i-1} \wedge u_i = \mu_i \\
& u \wedge u_i = \alpha_i
\end{aligned}
\right.
\]
\note{une rature après $\mu_i$}

\[
\left\lbrace
\begin{aligned}
& \lambda_i \mu_i + \lambda_{i+1} \mu_{i+1} = \alpha_i \\
& \lambda_i \alpha_{i-1} - \lambda_{i+1} \alpha_{i+1} = 0 \\
& \textstyle\sum \alpha_i = 0
\end{aligned}
\right.
\qquad\qquad 2 \sum \lambda_i \mu_i = 0
\]
\struck{(\ill{} \ill{} \ill{})}

$u_0, u_1, u, (\lambda_i) \Rightarrow u_i$ $\forall\, i$

$\alpha_0, \alpha_1, (\lambda_i) \Rightarrow (u_i)$ mod \ill{} Aff \quad
si $\mu_1 \neq 0$

\[
u = \frac{\alpha_1}{\mu_1} u_0 - \frac{\alpha_0}{\mu_1} u_1
= \Bigl( \lambda_1 + \lambda_2 \frac{\mu_2}{\mu_1} \Bigr) u_0
- \Bigl( \lambda_0 \frac{\mu_0}{\mu_1} + \lambda_1 \Bigr) u_1
\]
\note{une ligne biffée au-dessus}

$(\lambda_i, \mu_i)_i$ mod $k^*$.

$CI_n^* = n$-polygones cycliques (mod Aff) \struck{\ill{}} qui ne
\uncertain{sont} \ill{} ! \quad \struck{\ill{}} \struck{\ill{}}
\ill{}
\[
CI_n^* \longrightarrow \mathbf{P}^{n-1} \times \mathbf{P}^{n-1}, \qquad
(\lambda_i) \quad (\mu_i)
\]
\note{la flèche est verticale}

\marginal{Colonne de droite, séparée par un trait : $(\mu_i)_i \neq 0$ et
$u \neq 0$ $\iff$ $(\alpha_i)_i \neq 0$, le tout souligné d'une accolade
« hyp. » ; au-dessus, des essais très raturés (« $u \neq 0$ »,
« $\lambda_i \neq 0$ », « $\Rightarrow (u_i)_i \neq$ »). Puis, souligné :
« $u_i \neq 0$, $\lambda_i \neq 0$ $\forall\, i$ » hyp. ;
« $(\lambda_i, \alpha_i)_i \neq 0$ » ; « $\mathbf{P}(k^n \times
\uncertain{L}^n)$ »}

\[
u_2 = -\frac{\mu_2}{\mu_1} u_0 + \frac{1}{\lambda_2}
\Bigl( \lambda_0 \frac{\mu_0}{\mu_1} + \lambda_1 \Bigr) u_1
\]
\[
u_{n-1} = A_{n-1}\bigl((\lambda)(\mu)\bigr)\, u_0 + B_{n-1}(\ ,\ )\, u_1
\]
\note{devant $u_2$, un facteur $\lambda_2$ biffé, et un terme biffé
après le signe $=$}

\page{14}
\note{p.~1 de sa pagination, numéro cerclé en haut à droite ; la suite
court jusqu'à la page 28 (p.~16 de sa pagination, lot 2). Figure : le
sommet $s_i$ entre les côtés $u_{i-1}$ et $u_i$}
\[
f_i = x_{i-1} y_i - y_{i-1} x_i
\]
\[
df_i = y_i\, dx_{i-1} - x_i\, dy_{i-1} - (y_{i-1}\, dx_i - x_{i-1}\, dy_i)
\]
\[
\sum \lambda_i\, df_i = \sum_i (-\lambda_i y_{i-1} + \lambda_{i+1} y_{i+1})\,
dx_i + \sum_i (\lambda_i x_{i-1} - \lambda_{i+1} x_{i+1})\, dy_i
= \mu \sum dx_i + \nu \sum dy_i
\]
\note{une ligne biffée avant la dernière égalité}
\[
\left.
\begin{aligned}
-\mu &= +\lambda_i y_{i-1} - \lambda_{i+1} y_{i+1} \\
\nu &= +\lambda_i x_{i-1} - \lambda_{i+1} x_{i+1}
\end{aligned}
\right\rbrace \quad 0 \leqslant i \leqslant n-1
\]

\textbf{NB} \emph{OPS} $x_i, y_i$ tous $\neq 0$ (par \ill{} de
\uncertain{coord.})

\[
\lambda_0, \lambda_1 \ \big|\ \overset{i=1}{\lambda_2} \cdots
\overset{i=n-2}{\lambda_{n-1}} \qquad\qquad
\overset{i=n-1}{\lambda_n} \ \text{relation}
\]
\note{$\lambda_n$ est cerclé}

\[
\begin{aligned}
&+\lambda_0 y_{n-1} - \lambda_1 y_1 = +\lambda_1 y_0 - \lambda_2 y_2
= +\lambda_2 y_1 - \lambda_3 y_3 = \cdots
= \lambda_{n-2} y_{n-3} - \lambda_{n-1} y_{n-1} \\
&\qquad = \lambda_3 y_2 - \lambda_4 y_4 \qquad\qquad
= \lambda_{n-1} y_{n-2} - \lambda_n y_n \quad
(\lambda_n = \lambda_0,\ y_n = y_0)
\end{aligned}
\]

\[
\begin{aligned}
(y_2)\,\lambda_2 &= \lambda_1 (y_0 + y_1) - \lambda_0 y_{n-1}
 = \lambda_0 (-y_{n-1}) + \lambda_1 (y_0 + y_1) \qquad
 (= \lambda_4 y_3 - \lambda_5 y_5) \\
(y_2 y_3)\,\lambda_3 &= \underbrace{(y_2 \lambda_2)}\, y_1
 - y_2 (\lambda_0 y_{n-1} - \lambda_1 y_1) \\
 &= \lambda_0 \bigl( -y_{n-1} (y_1 + y_2) \bigr)
 + \lambda_1 y_1 [y_0 + y_1 + y_2] \\
(y_2 y_3 y_4)\,\lambda_4 &= (y_2 y_3 \lambda_3)\, y_3
 - y_2 y_3 (\lambda_0 y_{n-1} - \lambda_1 y_1) \\
 &= \lambda_0 \bigl( -y_2 y_{n-1} (y_1 + y_2 + y_3) \bigr)
 + \lambda_1 y_1 y_2 (y_0 + y_1 + y_2 + y_3) \\
(y_2 y_3 y_4 y_5)\,\lambda_5 &= \underbrace{(y_2 y_3 y_4 \lambda_4)}\, y_3
 - y_2 y_3 y_4 (\lambda_0 y_{n-1} - \lambda_1 y_1)
\end{aligned}
\]
\note{plusieurs indices sont surchargés dans ce calcul (le facteur après
$(y_2 y_3 \lambda_3)$ est écrit sur une rature) ; les formules sont données
telles qu'elles se lisent, sans les vérifier}

\page{15}
\note{p.~2 de sa pagination}
\[
= \lambda_0 \bigl( -y_{n-1} y_2 y_3 (y_1 + y_2 + y_3 + y_4) \bigr)
+ \lambda_1 y_1 y_2 y_3 (y_0 + y_1 + y_2 + y_3 + y_4)
\]
\note{puis une ligne biffée : « $(y_2 y_3 \cdots y_i)\lambda_i =
\lambda_0(-y_{n-1} y_2 y_3 \cdots y_{i-2}(y_1 + \cdots + y_{i-1})) +
\lambda_1 \ldots$ »}

\[
\boxed{
\begin{aligned}
(y_{i-1} y_i)\,\lambda_i &= (-\lambda_0 y_{-1})(y_1 + \cdots + y_{i-1})
 + (\lambda_1 y_1)(y_0 + \cdots + y_{i-1}) \\
(x_{i-1} x_i)\,\lambda_i &= (-\lambda_0 x_{-1})(x_1 + \cdots + x_{i-1})
 + (\lambda_1 x_1)(x_0 + \cdots + x_{i-1})
\end{aligned}}
\]
$i = 2, \ldots, n-1$ \quad $/\!/\ n$

\[
\boxed{\lambda_i u_{i-1} u_i = -\lambda_0 u_{-1} u_{1i} + \lambda_1 u_1 u_{0i}}
\]
où $u_{\alpha\beta} = u_\alpha + u_{\alpha+1} + \cdots + u_{\beta-1} =
s_\beta - s_\alpha$.

\[
\boxed{
\begin{aligned}
&\lambda_0 \bigl[ -x_{i-1} x_i y_{-1} (y_1 + \cdots + y_{i-1})
 + y_{i-1} y_i x_{-1} (x_1 + \cdots + x_{i-1}) \bigr] \\
&\quad + \lambda_1 \bigl[ x_{i-1} x_i y_1 (y_0 + \cdots + y_{i-1})
 - y_{i-1} y_i x_1 (x_0 + \cdots + x_{i-1}) \bigr] = 0
\end{aligned}}
\qquad i = 2, \ldots, n-1
\]
\note{« $i = 2, \ldots, n-1$ » est cerclé}

\[
(\underbrace{y_{n-1}}_{y_{-1}} \underbrace{y_n}_{y_0})
\underbrace{\lambda_n}_{\lambda_0}
= (-\lambda_0 y_{-1}) \underbrace{(y_1 + \cdots + y_{n-1})}_{-y_0}
+ (\lambda_1 y_1) \underbrace{(y_0 + \cdots + y_{n-1})}_{0}
\]
\[
\lambda_0 (y_{-1} y_0) = \lambda_0 y_0 y_{-1} \qquad \text{ok}
\]

\uncertain{triangle} \struck{\ill{}} :
\[
\begin{aligned}
&\lambda_0 \bigl( -x_1 x_2 y_{-1} \underbrace{(y_1)}_{-y_0}
 + y_1 y_2 x_{-1} (x_1) \bigr) \\
&\quad + x_1 x_2 y_2 (y_0) - y_1 y_2 x_2 (\ldots) \\
&\quad x_2 y_2 (x_1 y_0 - y_1 x_0)
\end{aligned}
\]
\note{calcul très raturé : les termes $y_2$, $x_2$ des parenthèses de la
première ligne sont biffés, de même que le contenu de la dernière
parenthèse de la deuxième ligne et une ligne entière
« $x_2 y_2 x_1 y_1 [-\frac{y_1 + y_2}{y_1} + \frac{x_1 + x_2}{x_1}$ » ;
le « $y_2$ » de « $x_1 x_2 y_2$ » est \uncertain{lu}}

\marginal{À droite, séparé par un trait oblique : $u$, $v$,
$w = -(u + v)$ ; $u \wedge v$, $v \wedge w$, $w \wedge u$ ;
$u \wedge v$, $-v \wedge u = u \wedge v$ \ill{}}

\page{16}
\note{p.~3 de sa pagination ; en tête, « \uncertain{quadrilatère} »}
\[
-x_1 x_2 y_1 y_3 + y_1 y_2 x_1 x_3 = -x_1 y_1 (x_2 y_3 - y_2 x_3)
\]
\note{le résultat est cerclé ; l'indice de $y_3$ est surchargé}

\[
\begin{aligned}
x_1 x_2 y_1 \underbrace{(y_0 + y_1)}_{-y_2 - y_3}
 - y_1 y_2 x_1 \underbrace{(x_0 + x_1)}_{-x_2 - x_3} &= \\
= -x_1 x_2 y_1 (y_2 + y_3) + y_1 y_2 x_1 (x_2 + x_3) & \\
= -x_1 x_2 y_1 y_3 + y_1 y_2 x_1 x_3 &= -x_1 y_1 (x_2 y_3 - y_2 x_3)
\end{aligned}
\qquad\qquad \boxed{\lambda_0 + \lambda_1 = 0}
\]
\note{le dernier membre est cerclé. Sous la formule encadrée, cerclé :
« si $s_3$ non \uncertain{aplati} », avec un croquis de l'angle en $s_3$
entre $u_2$ et $u_3$}

\[
-x_2 x_3 y_3 (y_1 + y_2) + y_2 y_3 x_3 (x_1 + x_2)
= -x_3 y_3 (x_1 y_2 - y_1 x_2)
\]
\[
x_2 x_3 y_1 \underbrace{(y_0 + y_1 + y_2)}_{-y_3}
- y_2 y_3 x_1 \underbrace{(x_0 + x_1 + x_2)}_{-x_3}
= x_3 y_3 (x_1 y_2 - y_1 x_2)
\qquad\qquad \boxed{\lambda_0 + \lambda_1 = 0}
\]
\note{les deux résultats sont cerclés, le signe du second surchargé ; des
points sont marqués sous certains facteurs. Sous le second encadré,
cerclé : « si $s_1$ non \uncertain{aplati} », avec le croquis de l'angle
en $s_1$ entre $u_1$ et $u_2$}

$u_0 \quad u_1 \quad u_2 \quad u_3$
\[
u_0 \wedge u_1 \quad u_1 \wedge u_2 \quad u_2 \wedge u_3 \qquad
u_3 \wedge u_0 = -(u_0 + u_1 + u_2) \wedge u_0 = u_0 \wedge u_1 +
u_0 \wedge u_2
\]
\note{une première forme de la dernière égalité est biffée ; le premier
terme du résultat est écrit sur une rature. En bas, un croquis d'angle
et le tableau des coordonnées $x_0, y_0$ ; $x_1, y_1$ ; $x_2, y_2$ ;
$x_3, y_3$}

\page{17}
\note{p.~4 de sa pagination ; en tête, « \uncertain{quadrilatère} ».
Figure : $s_0$ et $s_2$ sur une droite, $s_1$ et $s_3$ sur une droite
qui la coupe}
\[
\begin{array}{ll}
(a - 1, 0) & \bigl( -(a-1), b-1 \bigr) = (-a+1, b-1) \quad u_0 \\
(0, b - 1) & \bigl( a+1, -(b-1) \bigr) = (a+1, -b+1) \quad u_1 \\
(a + 1, 0) & \bigl( -(a+1), b+1 \bigr) = (-a-1, b+1) \quad u_2 \\
(0, b + 1) & \bigl( a-1, -(b+1) \bigr) = (a-1, -b-1) \quad u_3
\end{array}
\]
\note{la première ligne de la colonne de gauche est surchargée}

\[
\begin{aligned}
(1-a)(1-b) + (1-b)(a+1) &= 2(1-b) = f_1 \\
(a+1)(b+1) + (a+1)(1-b) &= 2(a+1) = f_2 \\
(a+1)(b+1) + (b+1)(1-a) &= 2(b+1) = f_3 \\
(a-1)(b-1) - (b+1)(a-1) &= 2(1-a) = f_0
\end{aligned}
\]
\note{le signe de la dernière ligne est surchargé, et les indices de
$f_1$ et $f_0$ sont écrits sur une rature}

\[
f_1 + f_3 = f_2 + f_0 = v \quad \Big/ \quad
\begin{aligned}
f_3 - f_1 &= (b)\, v \\
f_2 - f_0 &= (a)\, v
\end{aligned}
\]
\note{$v$ est écrit sur un « 4 » ; le calcul donne en effet $4$, et
$f_3 - f_1 = 4b$, $f_2 - f_0 = 4a$}

\note{figure : le quadrilatère de côtés $u_0, u_1, u_2, u_3$ et ses deux
diagonales fléchées}
\[
u_0 \wedge u_1 + u_2 \wedge u_3 = u_1 \wedge u_2 + u_3 \wedge u_0
\overset{?}{=} (u_0 + u_1) \wedge (u_1 + u_2)
= u_0 \wedge u_1 + u_0 \wedge u_2 + u_1 \wedge u_2
\]
\[
u_3 = -(u_0 + u_1 + u_2)
\]
\[
\underbrace{u_0 \wedge u_1} + u_2 \wedge \underbrace{(u_0 + u_1)}
= \underbrace{u_1 \wedge u_2} + u_0 \wedge \underbrace{(u_1 + u_2)}
\]
\note{dans les deux parenthèses, un troisième terme ($u_2$, puis $u_0$)
est barré d'une croix}
\[
u_2 \wedge \underbrace{(u_3 + u_0 + u_1)}_{-u_2} = 0 \qquad \text{ok}
\]

\note{sous une ligne ondulée : croquis de deux vecteurs $u$, $v$, avec
« $u \wedge v = -(-v) \wedge (-u)$ » biffé, et un pentagone avec toutes
ses diagonales}

\page{18}
\note{p.~5 de sa pagination. Figure : un 9-gone étoilé non convexe,
sommets $s_0, \ldots, s_8$, côtés $u_0, \ldots, u_7$, avec les segments
de $s_0$ aux autres sommets}
\[
\boxed{\lambda_0 = 0} \Rightarrow u_1 u_{0i} = \lambda_i u_{i-1} u_i
\]

\[
\begin{array}{lll}
s_0 = (0, 0) & & \\
s_1 = (1, 0) & & \\
s_2 = (1, 1) & & (y_2 = 1) \\
s_3 = (\beta_1, \beta_1) & \beta_1 \neq 1, 0 & \\
s_4 = (\beta_1, y_4) & y_4 \neq \beta_1 & \\
s_5 = (\beta_2 \beta_1, \beta_2 y_4) & \beta_2 \neq 1, 0 & \\
s_6 = (\beta_2 \beta_1, y_6) & y_6 \neq \beta_2 y_4 & \\
s_7 = (\beta_3 \beta_2 \beta_1, \beta_3 y_6) & \beta_3 \neq 1, 0 & \\
s_8 = (\beta_3 \beta_2 \beta_1, y_8) & y_8 \neq \beta_3 y_6 & \\
{[s_9 = s_0]} & [\beta_4 = 0] &
\end{array}
\qquad
\begin{array}{l}
\beta_1\ \beta_2\ \beta_3 \\
y_4\ y_6\ y_8
\end{array}
\]
\note{au-dessus des deux dernières lignes, « $\beta_1, \beta_2, \ldots,
\beta_8$ » biffé}

Aucun des $\lambda_i$ nuls \struck{$\iff$ \uncertain{les} $s_i$ tous
distincts}

\[
\lambda_i u_{i-1} - \lambda_{i+1} u_i - \lambda_{i+1} u_{i+1}
+ \lambda_{i+2} u_{i+2} = 0
\]
\[
\begin{array}{ll}
\lambda_1 u_0 - \lambda_2 u_1 - \lambda_2 u_2 + \lambda_3 u_3 = 0 &
\lambda_3 u_3 = -\lambda_1 u_0 + \lambda_2 u_1 + \lambda_2 u_2 \\
\lambda_2 u_1 - \lambda_3 u_2 - \lambda_3 u_3 + \lambda_4 u_4 = 0 &
\lambda_4 u_4 = -\lambda_1 u_0 + \ldots \\
\lambda_3 u_2 - \lambda_4 u_3 - \lambda_4 u_4 + \lambda_5 u_5 = 0 &
\lambda_3 \lambda_5 u_5 = \\
\lambda_4 u_3 - \lambda_5 u_4 - \lambda_5 u_5 + \lambda_6 u_6 &
\end{array}
\]
\note{à droite, « $(\lambda_2 + \lambda_3)\, u_2$ » et un $\lambda_3$
biffé}

\[
\begin{aligned}
\lambda_3 \lambda_5 u_5 &= \lambda_3 \bigl[ -\lambda_1 u_0
 + (\lambda_2 + \lambda_3)\, u_2 \bigr]
 + \lambda_4 (-\lambda_1 u_0 + \lambda_2 u_1 + \lambda_2 u_2)
 - \lambda_3^2 u_2 \quad \cdots \\
 &= \bigl( -\lambda_1 (\lambda_3 + \lambda_4) \bigr) u_0
 + \lambda_2 \lambda_4 u_1 + \lambda_2 (\lambda_3 + \lambda_4)\, u_2
\end{aligned}
\]
\note{puis un court essai biffé ; en bas, « $u_0\ u_1$ \quad 4 »}

\page{19}
\note{p.~6 de sa pagination. Figure : un pentagone $s_0 s_1 s_2 s_3 s_4$
de côtés orientés $u_0, \ldots, u_4$}
\[
\begin{array}{ll}
s_0 = (0, 0) & u_0 = (1, 0) \\
s_1 = (1, 0) & u_1 = (x - 1, y) \\
{[s_4 = (0, f_0)]} & u_2 = (x' - x, y' - y) \\
s_2 = (x, y) & u_3 = (-x', f_0 - y') \\
s_3 = (x', y') & u_4 = (0, -f_0)
\end{array}
\qquad
\begin{aligned}
f_1 &= y \\
f_2 &= (x - 1)(y' - y) - y(x' - x) \\
 &= (xy' - yx') - (y' - y) \\
f_3 &= (x' - x)(f_0 - y') + x'(y' - y) \\
 &= (xy' - yx') + f_0 (x' - x) \\
f_4 &= f_0 x' \\
f_0 &= f_0
\end{aligned}
\]

\[
\left\lbrace
\begin{aligned}
y &= f_1 \\
x' &= f_4/f_0 \\
xy' - y' &= f_2 + f_1 f_4/f_0 - f_1 \\
xy' - f_0 x &= f_3 + f_1 f_4/f_0 - f_4
\end{aligned}
\right.
\qquad\qquad
f_0 x - y' = -f_1 + f_2 - f_3 + f_4
\]
\note{quelques lettres biffées ou surchargées dans ce système (un
« $y' =$ » biffé à droite)}

\[
\left\lbrace
\begin{aligned}
y &= f_1 \\
x' &= f_4/f_0 \\
y' &= f_0 x + X_1^4
\end{aligned}
\right.
\qquad
X_1^4 = f_1 - f_2 + f_3 - f_4 \qquad
X = f_0 - f_1 + f_2 - f_3 + f_4 = f_0 - X_1^4
\]
\[
x (f_0 x + X_1^4) - f_0 x - X_1^4 = f_2 + f_1 f_4/f_0 - f_1
\]
\[
f_0 x^2 + \underbrace{(X_1^4 - f_0)}_{-X} x
- (f_1 - f_2 + f_3 - f_4 + f_2 + f_1 f_4/f_0 - f_1)
\]
\note{dans la dernière parenthèse, les termes qui se détruisent sont
barrés}
\[
f_0 x^2 - X x + (f_4 - f_3 - f_1 f_4/f_0) \qquad\qquad X = 1 - X_1^4
\]
\[
x = \frac{1}{2 f_0} \Bigl( X \pm \sqrt{X^2 - 4 f_0 f_4 + 4 f_0 f_3
+ 4 f_1 f_4} \Bigr)
\]
\note{sous le radical, la fin est surchargée ; on donne la forme que
demande la ligne précédente, les signes n'étant pas tous lisibles}
\[
\begin{aligned}
X^2 = f_0^2 + f_1^2 + f_2^2 + f_3^2 + f_4^2
 &- 2 f_0 f_1 + 2 f_0 f_2 - 2 f_0 f_3 + 2 f_0 f_4 \\
 &- 2 f_1 f_2 + 2 f_1 f_3 - 2 f_1 f_4 \\
 &- 2 f_2 f_3 + 2 f_2 f_4 \\
 &- 2 f_3 f_4
\end{aligned}
\]
\[
\Delta = f_0^2 + f_1^2 + f_2^2 + f_3^2 + f_4^2 + \ldots
\]
\note{les dix doubles produits de $\Delta$ sont écrits en colonne et
encadrés un à un, leurs signes surchargés ; ils ne sont pas repris ici}
\[
\begin{aligned}
&= \textstyle\sum f_i^2 - 2 \sum f_i f_{i+1} + 2 \sum f_i f_{i+2} \\
&= \textstyle\sum f_i^2 + 2 \sum f_{i+1} (f_i - f_{i+1}) \\
&\phantom{=}\ (f_0 - f_1 + f_2 + f_3 - f_4)^2
 + 4 (f_1 f_3 + f_2 f_3 + f_2 f_4)
\end{aligned}
\]
\note{la deuxième ligne est \uncertain{lue}}

\[
u_3 \wedge u_4 - u_2 \wedge u_3 = u_3 \wedge u_4 + u_3 \wedge u_2
= u_3 \wedge (u_2 + u_4) \qquad
(u_4 \wedge u_0)\bigl(u_3 \wedge (u_2 + u_4)\bigr) -
\]
\note{ligne séparée du reste par un trait, en bas de page}

\page{20}
\note{« 8 », non cerclé, en haut à droite : p.~8 de sa pagination (la
p.~7 manque). Un long trait de crayon traverse la page en diagonale}
\[
\begin{aligned}
&x^2 + y^2 + z^2 + 2yz + 2zx - 2xy = \\
&\xi^2 + \eta^2 - 2\xi\eta + 2(\xi + \eta) + 1 \\
&\underbrace{(\xi - \eta)}_{U}{}^2 + 2\underbrace{(\xi + \eta)}_{V} + 1
\end{aligned}
\]
\note{une ligne biffée entre la première et la deuxième}

\note{figure : quatre côtés consécutifs $u_{i-2}, \ldots, u_{i+1}$
annotés des $\lambda_{i-2}, \ldots, \lambda_{i+2}$, avec des ratures}

\struck{$\lambda_{i+1} u_i - \lambda_{i+2}(u_{i+1} + u_{i+2}) + \lambda_{i+3}
u_{i+3} = 0$}
\[
\boxed{\lambda_{i-1} u_{i-2} - \lambda_i (u_{i-1} + u_i) + \lambda_{i+1}
u_{i+1} = 0}
\]

\note{figure : $s_0$ et les côtés $u_0, u_1, u_2$ fléchés}
\[
\underbrace{u_0 \quad u_1}_{\text{base}} \qquad u_2 = \xi u_0 + \eta u_1
\]
\[
\begin{aligned}
u_3 &= \frac{1}{\lambda_3} \bigl( \lambda_2 (u_1 + u_2) - \lambda_1 u_0 \bigr) \\
u_4 &= \frac{1}{\lambda_3 \lambda_4} (\quad \\
u_5 &= \frac{1}{\lambda_3 \lambda_4 \lambda_5} (\quad \text{------} \quad)
\end{aligned}
\]
\[
u_{n-1} = \frac{\lambda_{n-2}}{\lambda_3 \cdots \lambda_{n-1}}
\bigl( P_{n-1} u_0 + Q_{n-1} u_1 + R_{n-1} u_2 \bigr)
\]
\[
u_i = \frac{1}{\lambda_3 \lambda_4 \cdots \lambda_i}\, P_i(\lambda_1, \ldots,
\lambda_{i-1})\, u_0 + Q_i(\lambda_1, \ldots
\]
\note{la dernière ligne s'arrête là ; elle est reprise, sous une autre
forme, en tête de la page 21 (lot 2). En marge, trois croquis de
polygones : un pentagone $s_0, \ldots, s_4$, et un $n$-gone
$s_0, s_1, s_2, \ldots, s_{n-1}$ de côtés $u_0, \ldots, u_{n-1}$}

\end{document}
